% thesis-preamble.tex
% ============================================================================
% LANGUAGE SUPPORT
% ============================================================================
\usepackage[english]{babel}
\usepackage{csquotes}

% ============================================================================
% MATHEMATICS PACKAGES
% ============================================================================
\usepackage{amsmath}
\usepackage{amssymb}
\usepackage{mathtools}
    \renewcommand{\[}{\begin{equation}}
    \renewcommand{\]}{\end{equation}}
\usepackage{nicefrac}
\usepackage{bbm}
\usepackage{siunitx}    % For units and numbers
\sisetup{
    table-align-text-post = false,
    table-format = 3.0,
    table-number-alignment = center
}
\usepackage{xparse}     % For advanced command definitions
 % Workaround for siunitx/physics compatibility
\NewCommandCopy\siqty\qty 
\usepackage{physics}  % Note: Must be loaded after siunitx if both are used

% ============================================================================
% GRAPHICS AND FIGURES
% ============================================================================
\usepackage{geometry}
\usepackage{graphicx}
\usepackage{subcaption}
\usepackage{adjustbox}
\usepackage[dvipsnames]{xcolor}
\usepackage{tikz}
\usetikzlibrary{positioning,calc,arrows.meta,fit,shapes,backgrounds,
                shapes.symbols,decorations.markings}
\usepackage{tikzscale}
\makeatletter
\newcommand{\vast}{\bBigg@{3}}
\newcommand{\Vast}{\bBigg@{5}}
\makeatother
% ============================================================================
% TABLES
% ============================================================================
\usepackage{multirow}
\usepackage{longtable}
% Column type for math mode in tables
\newcolumntype{M}{>{\(\displaystyle}l<{\)}}
\newcolumntype{C}{>{\(\displaystyle}c<{\)}}
\usepackage{threeparttable}
\usepackage{threeparttablex}

% ============================================================================
% SPECIALIZED PACKAGES
% ============================================================================
\usepackage[ruled]{algorithm}
\usepackage{algpseudocodex}
\algrenewcommand\algorithmiccomment[1]{\hfill\textit{\textcolor{gray}{// #1}}}
\algrenewcommand{\alglinenumber}[1]{\small\textcolor{gray}{#1}}
\algrenewcommand\algorithmicrequire{\textbf{Input:}}
\algrenewcommand\algorithmicensure{\textbf{Output:}}
\setlength{\algomargin}{1em}
\renewcommand{\algospacing}{\setlength{\baselineskip}{2\baselineskip}}
\usepackage{listings}
\usepackage{pdflscape}  % For landscape pages

% ============================================================================
% UTILITIES
% ============================================================================
\usepackage{appendix}   % For appendix formatting
\usepackage{enumitem}   % Better list environments
\usepackage{comment}    % Block comments
\usepackage{url}        % URL formatting

% ============================================================================
% BIBLIOGRAPHY
% ============================================================================
\usepackage[
    backend=biber,
    style=numeric-comp,
    sorting=none,
    maxbibnames=99,
    giveninits=true
]{biblatex}
\addbibresource{references.bib}

% ============================================================================
% CROSS-REFERENCES (load after most other packages)
% ============================================================================
\usepackage[
    hidelinks,
    pdfusetitle,
    bookmarks=true,
    bookmarksnumbered=true,
    bookmarksopen=true,
    bookmarksopenlevel=1,
    breaklinks=true,
    pdfborder={0 0 0},
    colorlinks=false
]{hyperref}

\usepackage[capitalise,noabbrev]{cleveref}
\crefname{paragraph}{Section}{Sections}
\crefname{subparagraph}{Section}{Sections}
\crefname{phfthm@theorem}{Theorem}{Theorems}
\crefname{phfthm@definition}{Definition}{Definitions}

% ============================================================================
% CUSTOM MATHEMATICAL NOTATION
% ============================================================================
% Common operators
\DeclareMathOperator*{\argmax}{\arg\!\max}
\DeclareMathOperator*{\argmin}{\arg\!\min}
\DeclareMathOperator{\E}{\mathbb{E}}
\DeclareMathOperator{\Var}{Var}
\DeclareMathOperator{\Cov}{Cov}

% Number sets
\newcommand{\R}{\mathbb{R}}
\newcommand{\N}{\mathbb{N}}
\newcommand{\Z}{\mathbb{Z}}
\newcommand{\C}{\mathbb{C}}
\renewcommand{\P}{\mathbb{P}}
% Common shortcuts
\newcommand{\inv}{^{-1}}
\newcommand{\qt}{\widetilde{q}}
\newcommand{\ASL}{\operatorname{ASL}^{\pm}}